\section{Introduction}

Constrained quadratic programs and Euclidean projections are ubiquitous in engineering, statistics, and applied mathematics. In recent work, Vestini and Kempf \cite{Vestini2025} addressed a critical limitation of Dykstra's projection algorithm, specifically the phenomenon of stalling when projecting an initial point onto the intersection of polyhedral sets. By identifying the exact duration of the stalling period and developing a fast-forwarding modification, the authors provided a computationally inexpensive method to circumvent arbitrarily long intervals during which the primal iterates remain unchanged. However, whilst the theoretical advancement and the corresponding stall-averse algorithm demonstrate robust convergence guarantees, the current literature lacks a concrete, large-scale application that inherently generates the highly acute polyhedral cones necessary to fully showcase the superiority of this modified approach in practice. 

This paper addresses that gap by applying the fast-forwarding Dykstra algorithm to the domain of nonlinear ensemble filtering, specifically within the context of optimal transport. We focus on the approximation of the Knothe-Rosenblatt (KR) rearrangement, a triangular transport map used to transform complex, non-Gaussian probability distributions into a standard normal reference measure. The computation of this map via maximum likelihood estimation requires enforcing strict monotonicity constraints across a dense ensemble of sample points. These constraints naturally form an ill-conditioned polyhedral feasible set characterised by nearly parallel intersecting half-spaces.

The primary contribution of this paper is the formulation of the projection step within a projected gradient descent (PGD) scheme for KR optimal transport, presenting it as a highly practical and rigorous application for the modified Dykstra algorithm. By parameterising the transport map with probabilist's Hermite polynomials, we construct a strictly convex optimisation framework constrained by a geometry where vanilla projection methods inevitably succumb to catastrophic stalling, thereby providing an ideal testbed for the fast-forwarding technique.

The remainder of this paper is structured as follows. Section \ref{sec:background} outlines the relevant background and related work. Section \ref{sec:core} details the core mathematical formulation, progressing sequentially from the underlying optimal transport problem to the specific projected gradient descent framework incorporating Dykstra's algorithm. Section \ref{sec:example} provides a fully worked two-dimensional example inspired by aerospace tracking applications.